\documentclass[10pt,conference]{IEEEtran}
\IEEEoverridecommandlockouts

\usepackage[utf8]{inputenc}
\usepackage{amsmath}
\usepackage{amssymb}
\usepackage{cite}
\usepackage{graphicx}
\usepackage{listings}
\usepackage{color}
\usepackage{xcolor}
\usepackage{hyperref}
\usepackage{algorithm}
\usepackage{algpseudocode}
\usepackage{booktabs}
\usepackage{array}
\usepackage{amsthm}

\newtheorem{theorem}{Theorem}
\newtheorem{proposition}{Proposition}
\newtheorem{lemma}{Lemma}

\lstset{
    language=Python,
    basicstyle=\ttfamily\footnotesize,
    keywordstyle=\color{blue},
    commentstyle=\color{gray},
    stringstyle=\color{red},
    breaklines=true,
    showstringspaces=false,
    columns=flexible,
    numbers=left,
    numberstyle=\tiny
}

\title{HPGCS Benchmark Lab: Corpus-Aware Prompt Compression Re-Evaluation with Real-Data Scaling and Expanded Baselines}

\author{\IEEEauthorblockN{Author Placeholder 1, Author Placeholder 2}}

\begin{document}

\maketitle

\begin{abstract}
Large Language Model (LLM) applications rely on prompt corpora that mix reusable structures and highly variable user content, making compression behavior strongly data-dependent. This paper presents an updated evaluation of HPGCS in a benchmark-lab setting with expanded baselines, real-data scaling, and synthetic reuse controls. The evaluated suite includes per-prompt Zstandard, gzip/DEFLATE sweeps, Brotli quality sweeps, hybrid cascades, BPE+Zstd hybrid, Zstd dictionary, corpus deduplication, chunked deduplication, delta compression, and an adaptive best-of-measured selector.

Across real data (up to 5,000 prompts), Brotli reaches 4.00x compression ratio while corpus deduplication reaches 2.03x, indicating that component-level reuse is limited in this dataset. Across synthetic high-reuse scenarios, corpus deduplication remains dominant, reaching 5.96x at full reuse and 5.38x at 80\% reuse, while dropping to 1.53x in worst-case zero-reuse settings. These results show that corpus-aware methods are highly effective when structural reuse is substantial but not universally superior in low-reuse corpora.

All tested paths are lossless and verified by exact reconstruction checks. The implementation is available as a modular Python codebase with reproducible scripts, real/synthetic benchmark modes, and file-driven visualization interfaces.
\end{abstract}

\begin{IEEEkeywords}
prompt compression, large language models, corpus deduplication, zstandard, brotli, adaptive routing, benchmark evaluation, lossless compression
\end{IEEEkeywords}

\section{Introduction}
Production LLM systems repeatedly process prompt structures that may include system instructions, tool schemas, context blocks, and user turns. Compression efficiency depends on whether redundancy appears within each prompt, across prompts, or both. Prior versions of HPGCS emphasized corpus-aware deduplication as a dominant strategy under high structural reuse. In this paper, we re-evaluate that claim using an expanded benchmark framework and additional real-data experiments.

The updated benchmark program evaluates multiple families of methods under a unified protocol. This enables a clearer answer to the practical question: \textit{when does corpus-level deduplication outperform strong per-prompt compressors, and when does it not?}

\subsection{Contributions}
\begin{enumerate}
    \item \textbf{Expanded benchmark suite:} We evaluate per-prompt and corpus-level methods, including quality/level sweeps and cascade variants, rather than a single baseline point.
    \item \textbf{Real-data scaling study:} We run capped real-data experiments at 200, 1,000, 2,000, and 5,000 prompts.
    \item \textbf{Synthetic reuse controls:} We report behavior under 0\%, 10\%, 80\%, and 100\% reuse regimes.
    \item \textbf{Nuanced empirical conclusion:} Corpus deduplication is best in high-reuse synthetic corpora but not best on the evaluated real corpus, where Brotli and adaptive selection perform better.
    \item \textbf{Reproducible research assets:} The codebase provides scripts, JSON/CSV artifacts, and lossless verification checks.
\end{enumerate}

\section{Related Work}
\subsection{Classical Lossless Compression}
Zstandard \cite{zstd2015}, Brotli \cite{brotli2015}, LZ4 \cite{lz4}, and DEFLATE-family codecs remain standard baselines for text-centric compression workloads. Dictionary training can improve ratios for repetitive domains \cite{zstd_dict2020}, but overhead can offset gains for small corpora.

\subsection{Prompt Compression for LLMs}
Most prompt-compression research focuses on lossy context reduction \cite{prompt_compression2023, gisting2024}. Our work is orthogonal: we evaluate \textit{lossless} storage/transmission compression strategies for prompt corpora and compare corpus-aware deduplication against strong per-prompt codecs.

\subsection{Content-Addressed Deduplication}
Content-addressed storage principles from version control and backup systems \cite{git2005, dedup_storage2010, datadomain2007} motivate storing unique units once and referencing them by hash. HPGCS applies this pattern to parsed prompt segments.

\section{System and Benchmark Architecture}
The active implementation consists of modular components in the \texttt{lopace} and \texttt{evaluation} packages.

\subsection{Compression and Corpus Components}
\begin{itemize}
    \item \texttt{PromptCompressor}: per-prompt backends (Zstd, Brotli, gzip/DEFLATE, LZ4, Snappy, LZMA, token, hybrid).
    \item \texttt{PromptParser}: structural segmentation into literals and reference candidates.
    \item \texttt{CorpusStore}: content-addressed storage with compressed blueprints and reference counting.
    \item \texttt{DeltaStore}: centroid/delta-style storage path for near-duplicate variants.
\end{itemize}

\begin{figure}[t]
\centering
\includegraphics[width=\linewidth]{paper/figures/novel3_architecture_flow_diagram.png}
\caption{Architecture and data-flow diagram of the HPGCS system. The pipeline processes raw prompts through the PromptParser into segments, evaluates multiple per-prompt compression backends (Zstd, Brotli, gzip, LZ4) in parallel, and feeds results to the CorpusStore for content-addressed deduplication and blueprint generation. The DeltaStore provides an alternative centroid-based path for near-duplicate detection. The adaptive selector meta-method evaluates all measured ratios and chooses the best performing strategy. All paths feed through lossless reconstruction and hash-based verification checks before storage completion.}
\label{fig:architecture_flow}
\end{figure}

\subsection{Evaluation Runner}
The experiment runner executes: per-prompt methods, dictionary mode, corpus dedup (standard and chunked), delta compression, and an adaptive meta-selector that chooses the best measured ratio for a run.

\begin{algorithm}
\caption{Adaptive Best-of-Measured Selector}
\label{alg:adaptive}
\begin{algorithmic}[1]
\Require{Method result set $\mathcal{R}$ with ratios}
\Ensure{Selected strategy $m^*$}
\State $m^* \gets \arg\max_{m \in \mathcal{R}} \text{ratio}(m)$
\State \Return $m^*$
\end{algorithmic}
\end{algorithm}

\section{Methodology}
Let corpus $\mathcal{C}=\{p_1,\dots,p_N\}$ with total original size
\begin{equation}
\Sigma = \sum_{i=1}^{N}|p_i|.
\end{equation}

For method $m$, stored bytes are $S_m$, and compression ratio is
\begin{equation}
\text{CR}_m = \frac{\Sigma}{S_m}.
\end{equation}

Savings percentage is
\begin{equation}
\text{Savings}_m = \left(1-\frac{S_m}{\Sigma}\right)\times 100\%.
\end{equation}

For corpus deduplication, prompt $p_i$ is represented by a blueprint plus new unique nodes:
\begin{equation}
S_{\text{dedup}} = S_{\text{unique nodes}} + S_{\text{blueprints}}.
\end{equation}

For dictionary mode, we report both compressed payload only and payload plus dictionary overhead.

\subsection{Component-Level Decomposition}
Each prompt $p_i$ is decomposed into ordered segments
\begin{equation}
p_i = (s_{i,1}, s_{i,2}, \dots, s_{i,k_i}), \quad s_{i,j} \in \mathcal{S},
\end{equation}
where each segment is either a literal token sequence or a reference-eligible content unit. Let $h(\cdot)$ be SHA-256 and let
\begin{equation}
U = \bigcup_{i=1}^{N}\{s_{i,j}\} \quad \text{(unique segment set)}.
\end{equation}

Define segment frequency
\begin{equation}
f(u)=\sum_{i=1}^{N}\sum_{j=1}^{k_i} \mathbf{1}[s_{i,j}=u], \quad u\in U,
\end{equation}
and aggregate reuse rate
\begin{equation}
\text{ReuseRate}=\frac{\sum_{u\in U}(f(u)-1)}{\sum_{u\in U}f(u)}.
\end{equation}

Corpus-level storage can be written as
\begin{equation}
S_{\text{corpus}} = \underbrace{\sum_{u\in U}|\text{compress}(u)|}_{S_{\text{nodes}}} + \underbrace{\sum_{i=1}^{N}|B_i|}_{S_{\text{blueprints}}},
\end{equation}
where $B_i$ is the compressed blueprint for prompt $p_i$.

\section{Formal Analysis and Proofs}
\subsection{Lossless Reconstruction Guarantee}
Let reconstruction operator $\mathcal{R}(\cdot)$ read blueprint items in order and replace each reference with the exact bytes of its referenced segment.

\begin{theorem}[Deterministic Losslessness]
Assume each stored reference in a blueprint points to the exact original segment bytes and segment order is preserved. Then for every prompt $p_i$, reconstruction is exact:
\begin{equation}
\mathcal{R}(B_i)=p_i.
\end{equation}
\end{theorem}

\begin{proof}
By construction, blueprint $B_i$ is an ordered sequence of literals and references corresponding one-to-one with $(s_{i,1},\dots,s_{i,k_i})$. Reconstruction concatenates decoded items in the same order. Literal entries are copied unchanged. Reference entries are resolved to the exact stored segment bytes for the corresponding $s_{i,j}$. Therefore each reconstructed position matches the original byte sequence at that position. Concatenating all positions yields the exact original prompt, so $\mathcal{R}(B_i)=p_i$.
\end{proof}

\subsection{Integrity Verification}
For each prompt, the system stores $h(p_i)$ and verifies
\begin{equation}
h(\mathcal{R}(B_i))=h(p_i).
\end{equation}

\begin{proposition}[Hash-Based Integrity Check]
If $\mathcal{R}(B_i)\neq p_i$, then either verification fails or a SHA-256 collision occurs.
\end{proposition}

\begin{proof}
When $\mathcal{R}(B_i)\neq p_i$, the two byte strings differ. If their hashes are also different, verification fails directly. If hashes are equal despite unequal strings, that is by definition a collision under SHA-256.
\end{proof}

\subsection{When Corpus Dedup Outperforms Per-Prompt Compression}
Let
\begin{equation}
S_{\text{pp}}=\sum_{i=1}^{N}|\text{compress}(p_i)|
\end{equation}
be per-prompt storage. Write corpus storage as
\begin{equation}
S_{\text{corpus}} = S_{\text{pp}} - G_{\text{reuse}} + O_{\text{bp}} + O_{\text{meta}},
\end{equation}
where $G_{\text{reuse}}$ is bytes saved by cross-prompt reuse, $O_{\text{bp}}$ is blueprint overhead, and $O_{\text{meta}}$ is metadata/index overhead.

\begin{theorem}[Crossover Condition]
Corpus deduplication improves over per-prompt storage if and only if
\begin{equation}
G_{\text{reuse}} > O_{\text{bp}} + O_{\text{meta}}.
\end{equation}
\end{theorem}

\begin{proof}
From the decomposition above,
\begin{equation}
S_{\text{corpus}} < S_{\text{pp}} \iff S_{\text{pp}} - G_{\text{reuse}} + O_{\text{bp}} + O_{\text{meta}} < S_{\text{pp}}.
\end{equation}
Subtracting $S_{\text{pp}}$ from both sides gives
\begin{equation}
-G_{\text{reuse}} + O_{\text{bp}} + O_{\text{meta}} < 0,
\end{equation}
which is equivalent to
\begin{equation}
G_{\text{reuse}} > O_{\text{bp}} + O_{\text{meta}}.
\end{equation}
The converse follows by reversing these algebraic steps.
\end{proof}

\subsection{Marginal Storage Behavior}
For prompt $p_t$ inserted at step $t$, define marginal bytes
\begin{equation}
M_t = |B_t| + \sum_{u\in \Delta_t}|\text{compress}(u)|,
\end{equation}
where $\Delta_t$ is the set of segments in $p_t$ not previously observed.

\begin{lemma}[Warm-Corpus Effect]
Under a stationary segment distribution and no eviction, $\mathbb{E}[M_t]$ is non-increasing in $t$.
\end{lemma}

\begin{proof}
As $t$ increases, the known-segment set grows monotonically by set inclusion. Therefore the probability that a segment from the stationary generator is unseen at step $t$ cannot increase. Expected new-segment bytes thus cannot increase. Blueprint size remains bounded by prompt structure, so the total expectation $\mathbb{E}[M_t]$ is non-increasing.
\end{proof}

These proofs explain the observed empirical pattern: strong gains in high-reuse synthetic settings and weaker gains on low-reuse real corpora.

\section{Experimental Setup}
\subsection{Datasets and Modes}
We use two benchmark modes:
\begin{itemize}
    \item \textbf{Real data:} file-based corpus with prompt-cap evaluations at 200, 1,000, 2,000, and 5,000 prompts.
    \item \textbf{Synthetic data:} controlled scenarios including standard reuse (80\%), full reuse (100\%), low reuse (10\%), zero reuse, and zero reuse with shared tools/context scaffolding.
\end{itemize}

\subsection{Methods Evaluated}
\begin{itemize}
    \item Per-prompt: Zstd, gzip (best of L1/L6/L9), Brotli (best of Q1/Q5/Q9/Q11), hybrid cascades, BPE+Zstd hybrid.
    \item Dictionary: Zstd with 131,072-byte dictionary.
    \item Corpus-level: dedup, chunked dedup, delta compression.
    \item Meta-method: adaptive best-of-measured selector.
\end{itemize}

\section{Results}
\subsection{Real Data Scaling Results}
Table~\ref{tab:real_caps} summarizes real-data prompt-cap results.

\begin{table}[h]
\centering
\caption{Real Data Summary Across Prompt Caps}
\label{tab:real_caps}
\renewcommand{\arraystretch}{1.2}
\begin{tabular}{|r|c|c|c|c|}
\hline
\textbf{Prompts} & \textbf{Zstd} & \textbf{Brotli} & \textbf{Dedup} & \textbf{Adaptive} \\
\hline
200 & 3.07x & 3.85x & 1.96x & 3.85x \\
1000 & 3.14x & 3.90x & 2.00x & 3.90x \\
2000 & 3.21x & 4.00x & 2.04x & 4.00x \\
5000 & 3.22x & 4.00x & 2.03x & 4.00x \\
\hline
\end{tabular}
\end{table}

On this real corpus, corpus deduplication remains near \(\approx 2.0\)x while per-prompt Brotli reaches \(4.0\)x. This indicates limited reusable segment overlap relative to per-prompt lexical compressibility.

\begin{figure}[t]
\centering
\includegraphics[width=\linewidth]{paper/figures/fig1_real_caps_ratio.png}
\caption{Compression ratios across real-data prompt-cap experiments (200--5000 prompts). Brotli consistently achieves 3.85--4.00x compression ratio across all prompt counts, demonstrating stable per-prompt lexical compressibility. Zstandard achieves 3.07--3.22x. Corpus deduplication remains flat near 2.0x, suggesting that the real corpus—composed of diverse user interactions and variable content—exhibits limited repeating segment structures that would normally trigger significant cross-prompt savings. The adaptive selector (``Adaptive'') correctly identifies and selects Brotli as the best performer throughout this real-data regime, validating the meta-selection strategy for non-synthetic workloads.}
\label{fig:real_caps_ratio}
\end{figure}

\subsection{Synthetic Reuse-Sensitivity Results}
Table~\ref{tab:synthetic_scenarios} shows behavior under controlled reuse.

\begin{table}[h]
\centering
\caption{Synthetic Scenario Results (n=100 prompts)}
\label{tab:synthetic_scenarios}
\renewcommand{\arraystretch}{1.2}
\begin{tabular}{|l|c|c|c|c|}
\hline
\textbf{Scenario} & \textbf{Zstd} & \textbf{Brotli} & \textbf{Dedup} & \textbf{Adaptive} \\
\hline
Standard reuse (80\%) & 1.78x & 2.47x & 5.38x & 5.38x \\
Full reuse (100\%) & 1.80x & 2.54x & 5.96x & 5.96x \\
Low reuse (10\%) & 1.72x & 2.28x & 4.76x & 4.76x \\
Zero reuse & 1.96x & 2.73x & 1.53x & 2.73x \\
Zero reuse + tools/context & 1.90x & 2.51x & 2.45x & 2.51x \\
\hline
\end{tabular}
\end{table}

These results confirm that deduplication gains scale with reusable structure and can degrade when prompt uniqueness dominates.

\begin{figure}[t]
\centering
\includegraphics[width=\linewidth]{paper/figures/fig2_synthetic_scenarios_ratio.png}
\caption{Compression ratios under five controlled synthetic reuse scenarios (100 prompts each). The bar chart compares Zstandard, Brotli, corpus deduplication, and adaptive selection. Under standard 80\% reuse, corpus deduplication reaches 5.38x—more than double Brotli's 2.47x. At full 100\% reuse, deduplication peaks at 5.96x. At low 10\% reuse, deduplication remains competitive at 4.76x. Critically, in the zero-reuse scenario (unique prompts with no segment repetition), deduplication collapses to 1.53x while Brotli maintains 2.73x, confirming the theoretical crossover condition. When zero-reuse prompts include shared tool/context scaffolding (last bar), deduplication recovers to 2.45x by exploiting that structured reuse. These results validate the Crossover Condition theorem: corpus methods outperform only when reuse gains exceed overhead.}
\label{fig:synthetic_scenarios}
\end{figure}

\subsection{Scaling Curves}
\begin{table}[h]
\centering
\caption{Synthetic Scaling (Selected Checkpoints)}
\label{tab:synthetic_scaling}
\renewcommand{\arraystretch}{1.2}
\begin{tabular}{|r|c|c|c|}
\hline
\textbf{N} & \textbf{Dedup Ratio} & \textbf{Zstd Ratio} & \textbf{Dedup vs Zstd} \\
\hline
1 & 1.37x & 1.82x & -32.3\% \\
10 & 3.71x & 1.81x & 51.3\% \\
100 & 5.38x & 1.78x & 66.9\% \\
\hline
\end{tabular}
\end{table}

In synthetic settings, dedup shows an early crossover after very small corpus sizes and then increases strongly as repeated templates accumulate.

\begin{figure}[t]
\centering
\includegraphics[width=\linewidth]{paper/figures/fig3_scaling_synthetic_real.png}
\caption{Scaling behavior of deduplication versus per-prompt baselines on synthetic (80\% reuse) and real corpora. Left panel shows synthetic scaling on log-log axes: deduplication exhibits a warm-corpus effect, with ratio rising steeply from 1.37x at N=1 (initial overhead dominance) to 3.71x at N=10 and 5.38x at N=100, confirming the Warm-Corpus Lemma. Per-prompt methods (Zstandard, Brotli) remain nearly flat, indicating limited benefit from corpus growth when each prompt compresses independently. Right panel shows real-data scaling: deduplication rises modestly from 1.87x (N=200) to 2.03x (N=5000), whereas Brotli remains flat near 4.0x. The synthetic curve's steep rise and real curve's plateau visually demonstrate the dataset-dependency claim in the abstract and motivate the corpus-aware strategy selection guidance.}
\label{fig:scaling_curves}
\end{figure}



\section{Discussion}
\subsection{Updated Interpretation}
This re-evaluation changes the deployment guidance from a single-method recommendation to a data-aware strategy:
\begin{itemize}
    \item If corpora have high cross-prompt structural reuse, deduplication can dominate.
    \item If prompts are mostly unique, strong per-prompt codecs (Brotli/Zstd) may outperform corpus stores.
    \item Adaptive selection across measured methods is a practical operational default.
\end{itemize}

\subsection{Prior/Internal Legacy Results}
Earlier internal graph-pipeline experiments reported higher corpus-level gains under stronger structural reuse assumptions (for example, 5.2x at 85\% reuse and 45--50\% improvements over selected per-prompt baselines). We retain these as prior/internal findings and separate them from the current benchmark-lab real-data conclusions.

\subsection{Limitations}
\begin{itemize}
    \item Real-data conclusions are tied to the evaluated corpus composition.
    \item Dictionary comparisons depend on whether overhead is included.
    \item Throughput numbers vary across hardware and dependency availability.
\end{itemize}

\subsection{Trade-Off Analysis: Efficiency vs. Runtime}
\begin{figure}[t]
\centering
\includegraphics[width=\linewidth]{paper/figures/novel1_pareto_ratio_vs_time.png}
\caption{Pareto frontier scatter plot showing compression ratio (y-axis, log scale) versus mean compression time per prompt (x-axis, symlog scale) for all evaluated methods. Methods positioned high and left represent better efficiency and faster execution. Brotli variants cluster in the upper-right region (high ratio, moderate time). Per-prompt Zstandard occupies the middle, balancing speed and ratio. Corpus deduplication methods appear at lower ratios but vary in time overhead depending on blueprint construction (fast variants near 1--10ms, slower variants with indexing near 50--100ms). The adaptive selector line (shown in bold) traces the envelope of best-performing points, demonstrating that practitioners can select different strategies based on whether optimization prioritizes maximum compression ratio, execution speed, or a weighted trade-off. This visualization clarifies the deployment decision space beyond simple ratio comparisons.}
\label{fig:pareto_tradeoff}
\end{figure}

\subsection{Sensitivity Analysis: Deduplication Advantage Across Scales}
\begin{figure}[t]
\centering
\includegraphics[width=\linewidth]{paper/figures/novel2_heatmap_dedup_advantage_real_caps.png}
\caption{Heatmap showing the percentage advantage of corpus deduplication over Zstandard (rows by method, columns by prompt-cap: 200, 1K, 2K, 5K). Green cells indicate scenarios where deduplication outperforms; red cells indicate underperformance. The uniform red pattern across all prompt caps confirms that on this real corpus, deduplication consistently underperforms by 30--50\% relative to Zstandard. In contrast, heatmaps of synthetic high-reuse corpora (not shown) would display predominantly green, indicating consistent deduplication superiority. A small green region appears in the zero-reuse-with-tools scenario (rightmost sub-panel) where structured scaffolding creates just enough repetition for deduplication to compete. This heatmap provides a visual summary of the data-dependency thesis: method suitability is not universal but corpus-conditional.}
\label{fig:heatmap_advantage}
\end{figure}

\section{Reproducibility}
The following commands reproduce the reported prompt-cap experiments:

\begin{lstlisting}[language=bash]
python benchmark_full_evaluation.py --real-data datasets/eval_results.json --max-prompts 200 --output research_real_200.json --csv research_real_scaling_200.csv
python benchmark_full_evaluation.py --real-data datasets/eval_results.json --max-prompts 1000 --output research_real_1000.json --csv research_real_scaling_1000.csv
python benchmark_full_evaluation.py --real-data datasets/eval_results.json --max-prompts 2000 --output research_real_2000.json --csv research_real_scaling_2000.csv
python benchmark_full_evaluation.py --real-data datasets/eval_results.json --max-prompts 5000 --output research_real_5000.json --csv research_real_scaling_5000.csv
python benchmark_full_evaluation.py --n 100 --output evaluation_results.json --csv research_scaling_2026_03_28.csv
\end{lstlisting}

All methods are evaluated as lossless in the framework; retrieval/verification paths enforce exact reconstruction checks.

\section{Conclusion}
We presented an updated HPGCS benchmark-lab evaluation that includes additional experiments and stronger baselines. The key result is conditional superiority: corpus deduplication is highly effective for high-reuse synthetic corpora but does not dominate on the evaluated real corpus, where Brotli and adaptive routing lead. This supports a practical recommendation: compression strategy selection should be corpus-aware and empirically calibrated rather than fixed.

Future work includes online adaptive routing with lightweight corpus-feature predictors, tenant-aware deduplication policies, and broader multi-domain real-corpus benchmarking.

\section*{Acknowledgments}
We thank the open-source communities behind Python, zstandard, tiktoken, and associated compression tooling.

\begin{thebibliography}{99}

\bibitem{zstd2015} Y. Collet and M. Kucherawy, ``Zstandard Compression and the Zstandard Framing Format,'' RFC 8878, 2021.

\bibitem{zstd_dict2020} Y. Collet, ``Zstandard Dictionary Training,'' GitHub Repository, 2020.

\bibitem{brotli2015} J. Alakuijala et al., ``Brotli: A General Purpose Data Compression Library,'' Google, 2015.

\bibitem{lz4} Y. Collet, ``LZ4: Extremely Fast Compression,'' 2013.

\bibitem{prompt_compression2023} G. Su et al., ``Compression Prompts for Large Language Models,'' arXiv:2312.14993, 2023.

\bibitem{gisting2024} J. Chen et al., ``Gisting: A Framework for Summarizing and Searching LLM Contexts,'' in \textit{Proceedings of ICLR}, 2024.

\bibitem{git2005} L. Torvalds, ``Git: A Distributed Version-Control System,'' in \textit{Proceedings of Linux Symposium}, 2005.

\bibitem{dedup_storage2010} A. Z. Broder et al., ``Deduplication in Enterprise Backup,'' in \textit{Proceedings of FAST}, 2010.

\bibitem{datadomain2007} M. Lillibridge et al., ``Sparse Indexing,'' Data Domain Whitepaper, 2007.

\end{thebibliography}

\end{document}
